
\documentclass{article}
\usepackage{colortbl}
\usepackage{makecell}
\usepackage{multirow}
\usepackage{supertabular}

\begin{document}

\newcounter{utterance}

\centering \large Interaction Transcript for game `cladder', experiment `full\_v1.5\_default', episode 2677 with qwen3.5{-}27b.
\vspace{24pt}

{ \footnotesize  \setcounter{utterance}{1}
\setlength{\tabcolsep}{0pt}
\begin{supertabular}{c@{$\;$}|p{.15\linewidth}@{}p{.15\linewidth}p{.15\linewidth}p{.15\linewidth}p{.15\linewidth}p{.15\linewidth}}
    \# & \multicolumn{2}{c}{Player} && \multicolumn{2}{c}{Game Master} \\
    \hline

    \theutterance \stepcounter{utterance}  
    & & & \multicolumn{4}{p{0.6\linewidth}}{
        \cellcolor[rgb]{0.9,0.9,0.9}{
            \makecell[{{p{\linewidth}}}]{
                \texttt{\tiny{[P1$\langle$GM]}}
                \texttt{You are an expert in causal inference. The following question is not a typical commonsense query, but rather a meticulously designed question created by a professor specializing in causal inference, intended to assess the students' mastery of the course content.} \\
\texttt{Imagine a self{-}contained, hypothetical world with only the following conditions, and without any unmentioned factors or causal relationships: Gene has a direct effect on smoking and lung cancer. Pollution has a direct effect on smoking and lung cancer. Smoking has a direct effect on lung cancer. Pollution is unobserved. Method 1: We look at how gene correlates with lung cancer case by case according to smoking. Method 2: We look directly at how gene correlates with lung cancer in general. To understand how gene affects lung cancer, is it more correct to use the Method 1 than Method 2?} \\
\texttt{Start your answer with 'Yes' or 'No', followed by additional reasoning or evidence to support your explanation.} \\
            }
        }
    }
     \\ \\

    \theutterance \stepcounter{utterance}  
    & \multicolumn{4}{p{0.6\linewidth}}{
        \cellcolor[rgb]{0.561,0.737,0.561}{
            \makecell[{{p{\linewidth}}}]{
                \texttt{\tiny{[P1$\rangle$GM]}}
                \texttt{No} \\
\\ 
\texttt{In this hypothetical scenario, **Method 2** (looking at the direct correlation between Gene and Lung Cancer in general) is the correct approach to estimate the total causal effect of Gene on Lung Cancer, whereas **Method 1** (stratifying by Smoking) introduces bias.} \\
\\ 
\texttt{Here is the causal reasoning based on the structure provided:} \\
\\ 
\texttt{1.  **Causal Structure**:} \\
\texttt{    *   \$Gene \textbackslash{}rightarrow Smoking\$} \\
\texttt{    *   \$Gene \textbackslash{}rightarrow LungCancer\$} \\
\texttt{    *   \$Pollution \textbackslash{}rightarrow Smoking\$} \\
\texttt{    *   \$Pollution \textbackslash{}rightarrow LungCancer\$} \\
\texttt{    *   \$Smoking \textbackslash{}rightarrow LungCancer\$} \\
\texttt{    *   \$Pollution\$ is unobserved.} \\
\\ 
\texttt{2.  **Analysis of Method 2 (Marginal Association)**:} \\
\texttt{    *   This method estimates the association \$P(LungCancer | Gene)\$.} \\
\texttt{    *   The paths from Gene to Lung Cancer are:} \\
\texttt{        *   Direct path: \$Gene \textbackslash{}rightarrow LungCancer\$} \\
\texttt{        *   Indirect path: \$Gene \textbackslash{}rightarrow Smoking \textbackslash{}rightarrow LungCancer\$} \\
\texttt{    *   Since \$Pollution\$ is a common cause of \$Smoking\$ and \$LungCancer\$ but is **not** a cause of \$Gene\$, there is no backdoor path from \$Gene\$ to \$LungCancer\$ that needs to be blocked. The path \$Gene \textbackslash{}leftarrow \textbackslash{}dots \textbackslash{}rightarrow LungCancer\$ does not exist.} \\
\texttt{    *   Therefore, the marginal association in Method 2 captures the **Total Causal Effect** of Gene on Lung Cancer (both direct and indirect via smoking) without confounding bias.} \\
\\ 
\texttt{3.  **Analysis of Method 1 (Stratification by Smoking)**:} \\
\texttt{    *   This method attempts to estimate the effect by conditioning on \$Smoking\$ (i.e., looking at \$P(LungCancer | Gene, Smoking)\$).} \\
\texttt{    *   \$Smoking\$ is a **collider** on the path \$Gene \textbackslash{}rightarrow Smoking \textbackslash{}leftarrow Pollution\$.} \\
\texttt{    *   In causal inference, conditioning on a collider (or a descendant of a collider) opens a non{-}causal path between its parents. By stratifying on \$Smoking\$, we induce a spurious correlation between \$Gene\$ and \$Pollution\$.} \\
\texttt{    *   Once this path is opened, a backdoor path is created: \$Gene \textbackslash{}leftrightarrow Pollution \textbackslash{}rightarrow LungCancer\$.} \\
\texttt{    *   Since \$Pollution\$ is unobserved, we cannot adjust for it. Consequently, the estimate derived from Method 1 will be **biased** because it conflates the effect of Gene with the effect of the unobserved Pollution, which is now correlated with Gene due to the conditioning.} \\
\\ 
\texttt{**Conclusion**:} \\
\texttt{Method 1 creates "collider bias" (also known as Berkson's paradox in this context) by conditioning on a variable (\$Smoking\$) that is a common effect of the treatment (\$Gene\$) and an unobserved confounder (\$Pollution\$). Method 2 avoids this error and correctly estimates the total effect. Therefore, it is **not** more correct to use Method 1; Method 2 is the valid approach.} \\
            }
        }
    }
    & & \\ \\

    \theutterance \stepcounter{utterance}  
    & & & \multicolumn{2}{p{0.3\linewidth}}{
        \cellcolor[rgb]{0.9,0.9,0.9}{
            \makecell[{{p{\linewidth}}}]{
                \texttt{\tiny{[GM$|$GM]}}
                \texttt{No} \\
            }
        }
    }
    & & \\ \\

    \theutterance \stepcounter{utterance}  
    & & & \multicolumn{2}{p{0.3\linewidth}}{
        \cellcolor[rgb]{0.9,0.9,0.9}{
            \makecell[{{p{\linewidth}}}]{
                \texttt{\tiny{[GM$|$GM]}}
                \texttt{no} \\
            }
        }
    }
    & & \\ \\

    \theutterance \stepcounter{utterance}  
    & & & \multicolumn{2}{p{0.3\linewidth}}{
        \cellcolor[rgb]{0.9,0.9,0.9}{
            \makecell[{{p{\linewidth}}}]{
                \texttt{\tiny{[GM$|$GM]}}
                \texttt{game\_result = WIN} \\
            }
        }
    }
    & & \\ \\

\end{supertabular}
}

\end{document}
